\section{Discussion}

\subsection{Interpreting the Detection Signal}

The 8$\times$ loss difference provides a statistically robust signal. However, translating distributional difference into individual identification is constrained by base rate mathematics: at 5\% minority rate, perfect ranking achieves $\sim$33\% precision---which we attain.

The appropriate use case is weighted training (as in JTT) rather than hard minority labeling.

\subsection{Magnitude vs.\ Timing Gap}

With ImageNet pretraining, both feature types peak at epoch 81. Simplicity bias manifests as consistently \emph{higher} spurious encoding rather than \emph{earlier} encoding. Training from scratch may show different dynamics.

\subsection{Limitations}

\begin{itemize}
    \item \textbf{Single dataset}: Only Waterbirds tested
    \item \textbf{Pretrained only}: From-scratch dynamics may differ
    \item \textbf{No intervention evaluation}: Detection characterized, not applied
\end{itemize}
